% !TeX root = ../thuthesis-example.tex

\chapter{研究方法}

本章给出本文采用的研究方法。与早期将路径质量写成“链路保真度乘积乘以统一指数衰减”的模型不同，当前工作直接以 qns-3 中已经实现的路径标签语义为准：每条候选路径保存完整链路段列表，扩展时重新计算 $T_{\max}$、等待退相干和 Bell 对角态交换。

\section{路径级保真度模型}

对候选路径
\[
P=(v_0,v_1,\ldots,v_h),
\]
第 $j$ 段链路 $e_j=(v_{j-1},v_j)$ 在路由 label 中保存初始保真度、噪声族、链路就绪时间和两端节点相干时间。路由计算首先取
\[
T_{\max}(P)=\max_j T_j,
\]
然后令第 $j$ 段链路等待
\[
W_j(P)=T_{\max}(P)-T_j.
\]
等待过程由单量子比特相位翻转存储信道给出，等待后的 Bell 对角态再按路径顺序递归做纠缠交换。最终 Bell 对角态相对于 $\vert\Phi^+\rangle$ 的权重即为 \texttt{predicted\_fidelity}。

该模型有三个重要后果：第一，路径质量不是可加边权和；第二，新增一条边可能改变整条前缀的 $T_{\max}$，从而改变前面所有链路的等待时间；第三，同一节点处局部更优的前缀在拼接同一后缀后不一定继续更优。

\section{算法设计}

本文将 \texttt{DijkstraRoutingProtocol} 作为单标签基线。它为每个节点只保留一个当前最优 label，并使用 \texttt{predicted\_fidelity} 作为主排序键，\texttt{t\_max\_ms} 和路径长度作为并列情况下的次级排序键。该算法可以看作当前非加性度量下的 $K=1$ best-first 搜索，而不是经典加性意义下具有最优性保证的 Dijkstra。

为缓解单标签剪枝导致的前缀陷阱，本文采用 \texttt{SlicedDijkstraRoutingProtocol}。它允许同一节点保留至多 $K$ 条 label，并可沿 $T_{\max}$ 轴按固定桶宽 \texttt{BucketWidthMs} 分桶：
\[
\texttt{bucket\_id}(P)
=
\left\lfloor \frac{T_{\max}(P)}{\texttt{BucketWidthMs}} \right\rfloor .
\]
启用 bucket 时，不同时间切片中的 label 可以共存；同一切片内则保留 \texttt{IsBetter} 意义下更优的代表。关闭 bucket 时，算法退化为基于 \texttt{predicted\_fidelity} 和 \texttt{t\_max\_ms} 的 Pareto 多标签搜索。

需要注意的是，切片算法不是全局路径枚举。有限 $K$ 与同桶竞争规则仍可能丢弃未来更优的前缀，因此本文把它定位为缓解非保序剪枝的多标签近似，而不是严格最优算法。

\section{理论分析方法}

本文采用构造性反例分析当前度量的非保序性。核心思路是构造两条到达同一中间节点的前缀 $a,b$，使得
\[
F_{\mathrm{pred}}(a)>F_{\mathrm{pred}}(b),
\qquad
T_{\max}(a)>T_{\max}(b),
\]
再接入同一快速后缀 $x$。由于 $a$ 携带更大的时间瓶颈，后缀 $x$ 在 $a\oplus x$ 中需要更久等待；该等待通过存储退相干降低后缀 Bell 对角态质量，并最终在交换后导致
\[
F_{\mathrm{pred}}(a\oplus x)
<
F_{\mathrm{pred}}(b\oplus x).
\]
该反例直接说明单标签 Dijkstra 的局部最优前缀不能安全代表所有后续扩展。

\section{实验设计思路}

当前论文版本不把实验数值作为重点。后续实验将分为三类：第一，在手工构造的前缀陷阱拓扑上验证单标签剪枝与切片保留的差异；第二，在随机拓扑上统计 $K$、\texttt{BucketWidthMs} 与路由结果之间的关系；第三，在真实拓扑上用混态仿真复核路由 metric 对最终端到端保真度的预测能力。由于当前实际混态重跑尚未完成，本文只保留实验设计，不在理论章节中提前给出结果性结论。
